% Part: sets-functions-relations
% Chapter: functions
% Section: composition

\documentclass[../../../include/open-logic-section]{subfiles}

\begin{document}

\olfileid{sfr}{fun}{cmp}
\olsection{Samenstelling van functies}

\begin{explain}
\oliflabeldef{sfr:fun:inv:sec}{In \olref[inv]{sec} zagen we dat de
inverse~$f^{-1}$ van !!a{bijection}~$f$ zelf een functie is. Een andere
bewerking op functies is samenstelling: w}{W}e kunnen een nieuwe
functie definiëren door twee functies, $f$ en~$g$, samen te stellen:
we passen eerst $f$ toe en daarna~$g$. Dit kan alleen als hun beelden
en domeinen op elkaar aansluiten: het beeld van~$f$ moet een
deelverzameling zijn van het domein van~$g$.
\oliflabeldef{sfr:rel:ops:sec}{Deze bewerking met functies is vergelijkbaar
met de samenstelling van relaties uit
\olref[rel][ops]{sec}.}{}

Een diagram kan het idee van samenstelling verduidelijken. In
\olref{fig:composition} tonen we twee functies $f \colon A \to B$ en
$g \colon B \to C$, met hun samenstelling~$(\comp{f}{g})$. De functie
$(\comp{f}{g}) \colon A \to C$ koppelt elk !!{element} van~$A$ aan
!!a{element} van~$C$. Welk !!{element} van~$C$ bij !!a{element} van
$A$ hoort, leggen we als volgt vast. Gegeven een invoer $x \in A$
passen we eerst de functie $f$ toe op~$x$; dit levert $f(x) = y \in B$
op. Vervolgens passen we de functie $g$ toe op~$y$; dit levert
$g(f(x)) = g(y) = z \in C$ op.
\begin{figure}
  \olasset[2\olphotowidth]{assets/diagrams/composition.tikz}
  \caption{De samenstelling $g \circ f$ van twee functies $f$ en~$g$.}
  \ollabel{fig:composition}
\end{figure}
\end{explain}

\begin{defn}[Samenstelling]
Zij $f\colon A \to B$ en $g\colon B \to C$ functies. De
\emph{samenstelling} van $f$ met~$g$ is $\comp{f}{g} \colon A \to C$,
waarbij $(\comp{f}{g})(x) = g(f(x))$.
\end{defn}

\begin{ex}
Beschouw de functies $f(x) = x + 1$ en $g(x) = 2x$. Omdat
$(\comp{f}{g})(x) = g(f(x))$, nemen we voor elke invoer~$x$ eerst de
opvolger en vermenigvuldigen we het resultaat vervolgens met twee. Hun
samenstelling wordt dus gegeven door $(\comp{f}{g})(x) = 2(x+1)$.
\end{ex}

\begin{prob}
Toon aan dat als $f \colon A \to B$ en $g \colon B \to C$ beide
!!{injective} zijn, ook $\comp{f}{g}\colon A \to C$ !!{injective} is.
\end{prob}

\begin{prob}
Toon aan dat als $f \colon A \to B$ en $g \colon B \to C$ beide
!!{surjective} zijn, ook $\comp{f}{g}\colon A \to C$ !!{surjective}
is.
\end{prob}

\begin{prob}
Stel dat $f \colon A \to B$ en $g \colon B \to C$. Toon aan dat de
grafiek van $\comp{f}{g}$ gelijk is aan $R_f \mid R_g$.
\end{prob}
\end{document}
